\chapter{System Block Diagram}

\begin{figure}[h]
\centering
\begin{tikzpicture}[
  node distance=1.8cm and 1.2cm,
  block/.style={draw, rounded corners, minimum width=3.3cm, minimum height=1cm, align=center},
  line/.style={-Latex, thick}
]
\node[block] (client) {Remote Client};
\node[block, right=of client] (grpc) {gRPC Signing Daemon\\Linux on Zynq PS};
\node[block, right=of grpc] (axi) {AXI-Lite MMIO\\Wrapper Interface};
\node[block, right=of axi] (core) {Future ML-DSA-87\\Signing Core};
\node[block, below=of core] (keys) {PoC Key Storage\\in PL};

\draw[line] (client) -- node[above] {Ethernet} (grpc);
\draw[line] (grpc) -- node[above] {MMIO transactions} (axi);
\draw[line] (axi) -- node[above] {control, status, data} (core);
\draw[line] (keys) -- (core);
\end{tikzpicture}
\caption{Initial proof-of-concept system partition.}
\end{figure}

The diagram emphasizes the intended separation between the network service, PS-side orchestration, wrapper boundary, and PL-resident signing function.
